% !TEX encoding = UTF-8
% !TEX TS-program = pdflatex
% !TEX root = ../main.tex
% !TEX spellcheck = it-IT

\documentclass[final, 11pt, a4paper, titlepage]{article}
\makeatletter
\AtBeginDocument{\let\hl\@firstofone}
\makeatother

\usepackage[italian, english]{babel}
\usepackage[utf8]{inputenc}
\usepackage[hidelinks]{hyperref}
\usepackage{graphicx}
\usepackage{textcomp}
\usepackage{wallpaper}
\usepackage{listings}
\usepackage{color}
\usepackage{mathtools}
\usepackage{amssymb}
\usepackage{bussproofs}
\usepackage{listing}
\usepackage{bm}
\usepackage[
	backend=biber,
	citestyle=numeric-comp,
	hyperref,
	backref,
	sorting=none
]{biblatex}

\graphicspath{{images/}}

\addbibresource{bibliography.bib}


\defbibheading{bibliography}
{
    \phantomsection 
    \addcontentsline{toc}{section}{\bibname}
    \section*{\bibname\markboth{\bibname}{\bibname}}
}

\setlength\bibitemsep{1.5\itemsep} 

\DeclareBibliographyCategory{papers}

\addtocategory{papers}{Schwartz:2010:YEW:1849417.1849981}
\addtocategory{papers}{conf/sp/RochaBHWE10}
\addtocategory{papers}{Enck:2010:TIT:1924943.1924971}
\addtocategory{papers}{SlowinskaBos09}
\addtocategory{papers}{Sharma12acritical}


\newcommand{\university}{Università degli Studi di Padova}
\newcommand{\dept}{Department of Mathematics}
\newcommand{\faculty}{Master Degree in Computer Science}
\newcommand{\myyear}{Academic Year 2016-2017}
\renewcommand{\title}{Taint Mechanisms}
\newcommand{\subtitle}{Advanced Topics in Computer Security}
\renewcommand{\author}{Matteo Di Pirro}
\newcommand{\matr}{1154231}

\begin{document}
	\begin{titlepage}
\begin{center}

\begin{LARGE}
\textbf{\university}\\
\end{LARGE}

\vspace{10pt}

\begin{Large}
\textsc{\dept}\\
\end{Large}

\vspace{10pt}

\begin{large}
\textsc{\faculty}\\
\end{large}

\vspace{30pt}
\begin{figure}[htbp]
\begin{center}
\includegraphics[height=6cm]{images/logo_unipd}
\end{center}
\end{figure}
\vspace{30pt} 

\begin{LARGE}
\begin{center}
\textbf{\title}\\
\end{center}
\end{LARGE}

\begin{Large}
	\begin{center}
		\textbf{\subtitle}\\
	\end{center}
\end{Large}

\vspace{20pt} 

\begin{large}
\begin{flushright}
\textit{\author \\ \matr}
\end{flushright}
\end{large}

\vspace{40pt}

\line(1, 0){338} \\
\begin{normalsize}
\textsc{\myyear}
\end{normalsize}

\end{center}
\end{titlepage} 
	\tableofcontents
	\newpage
	\section{Introduction}
\subsection{Static and Dynamic Analysis} 
Dynamic analysis is the analysis of the properties of a running program. In contrast to static analysis, which examines a program's text to derive properties that hold for all executions, dynamic analysis derives properties that hold for one or more executions by examination of the running program. While dynamic analysis cannot prove that a program satisfes a particular property, it can detect violations of properties as well as provide useful information to programmers about the behavior of their programs. Although dynamic analysis provides a powerful mechanism for relating program inputs, this dipendence from the inputs makes it incomplete. Viewed in this light, dynamic and static analysis might be better termed ''input-centric'' and ''program-centric'' analysis respectively. \cite{ball1999concept}

Dynamic analysis is attractive because it allows us to reason about actual executions, and thus can perform precise security analysis based upon run-time information. Further, dynamic analysis is simple: we need only consider facts about a single execution at a time. 

\subsection{Questions about user input}
The two analyses can be used in conjunction to build formulas representing only the parts of an execution that depend upon tainted values.

\subsubsection{Is the final value affected by user input? $\Rightarrow$ Dynamic Taint Analysis}
Dynamic taint analysis runs a program and observes which computations are affected by predefined taint sources such as user input. . Any program value whose computation depends on data derived from a taint source is considered \textit{tainted}. Any other value is considered untainted.

\subsubsection{What input will make execution reach this line of code? $\Rightarrow$ Forward Symbolic Execution}
Dynamic forward symbolic execution automatically builds a logical formula describing a program execution path, which reduces the problem of reasoning about the execution to the domain of logic.

\subsection{Applications}
The number of security applications utilizing these two techniques is enormous. Example security research areas employing either dynamic taint analysis, forward symbolic execution, or a mix of the two, are te following:
\begin{itemize}
	\item \textbf{Unknown Vulnerability Detection}: Dynamic taint analysis can look for misuses of user input during an execution. For example, dynamic taint analysis can be used to prevent code injection attacks by monitoring whether user input is executed;
	\item \textbf{Automatic Input Filter Generation}: Forward symbolic 	execution can be used to automatically generate input filters that detect and remove exploits from the input stream. Filters generated in response to actual executions are attractive because they provide strong accuracy guarantees;
	\item \textbf{Malware Analysis}: Taint analysis and forward symbolic execution are used to analyze how information flows through a malware binary, explore trigger-based behavior and detect emulators;
	\item \textbf{Test Case Generation}: Taint analysis and forward symbolic execution are used to automatically generate inputs to test programs, and can generate inputs that cause two implementations of the same protocol to behave differently.
\end{itemize}

	\begin{frame}%[fragile]
	\frametitle{Dynamic Taint Analysis}
	{\fontsize{0.8cm}{0.5cm}\selectfont
		\texttt{x := get\_input(\includegraphics[scale=0.05]{devil}) \\ \mbox{} \\
		z := 42 \\ \mbox{} \\
		y := x + z\\ \mbox{} \\
		goto y\\}
	}
	\begin{tikzpicture}[remember picture,overlay]
	\node<2>[draw,minimum size=0.9cm,line width=1pt,white,circle,label=above:{\color{white}Tainted}] at (0.18,3.68) {};
	\node<2>[white,align=center,font=\fontsize{0.8cm}{0.5cm}\selectfont] at (7.2,2) {x is derived from\\a tainted source};
	
	\node<3>[draw,minimum size=0.9cm,line width=1pt,white,circle,label=above:{\color{white}Untainted}](z) at (0.18,2.62) {};
	\node<3>[white,align=center,font=\fontsize{0.8cm}{0.5cm}\selectfont] at (7.2,2) {z is a "static"\\constant};
	%\node<3>[white,font=\fontsize{0.6cm}{0.5cm}\selectfont](q) at (7.2,2.15) {Untainted};
	%\draw<3> [->,white,line width=1pt] (0.55,2.15) -- (q);
	
	\node<4->[draw,minimum size=0.9cm,line width=1pt,white,circle](y) at (0.18,1.6) {};
	\node<4->[white,font=\fontsize{0.6cm}{0.5cm}\selectfont](q) at (7.2,1.4) {Is y taited?};
	\draw<4-> [->,white,line width=1pt] (0.57,1.4) -- (q);
	
	\node<5>[white,align=center,font=\fontsize{0.8cm}{0.5cm}\selectfont] at (7.2,0) {It depends on the\\selected policy};
	\end{tikzpicture}
\end{frame}
	\begin{frame}[fragile]
	\frametitle{Forward Symbolic Execution}
	\begin{itemize}
		\item We can reason about the behavior of a program using the logic...
		\item ... and it is \textbf{conceptually} a very simple process
	\end{itemize}
	\lstset{language=C++,
		basicstyle=\color{white}\ttfamily,
		commentstyle=\color{green}\ttfamily,
		%		numbers=left
	}
	\begin{lstlisting}
	x := 2 * get_input(src)
	if x - 5 == 14 then goto 3 else goto 4
	// line 3: catastrophic failure
	// line 4: normal behaviour
	\end{lstlisting}
	\only<2->{
		\begin{itemize}
			\item \texttt{get\_input(src)} now returns a \textbf{symbol} instead of a concrete value
			\item But now expressions \textbf{cannot} be fully evaluated to a concrete value
		\end{itemize}
	}
\end{frame}
	\section{Conclusions}
\subsection{Comparison}
In my opinion, symbolic execution uses a kind of taint analysis to construct path constraints. Symbolic execution also employs an SMT solver to generate concrete values for variables and/or inputs, such that a certain path constraint is satisfied. Since taint analysis does not employ an SMT solver, I would say it is not a kind of symbolic execution.

Furthermore, dynamic analysis can also reason about \textit{feasible} path, while symbolic execution builds constraint and reasons about the whole program by reducing it in the domain of logic.

Thus, in my opinion, taint analysis is a \textbf{subset} of symbolic execution.

\subsection{Last things}
The number of security applications utilizing these two techniques is enormous. Example security research areas employing either dynamic taint analysis, forward symbolic execution, or a mix of the two, are the following:
\begin{itemize}
	\item \textbf{Unknown Vulnerability Detection}: Dynamic taint analysis can look for misuses of user input during an execution. For example, dynamic taint analysis can be used to prevent code injection attacks by monitoring whether user input is executed;
	\item \textbf{Malware Analysis}: Taint analysis and forward symbolic execution are used to analyze how information flows through a malware binary, explore trigger-based behavior and detect emulators;
	\item \textbf{Test Case Generation}: Taint analysis and forward symbolic execution are used to automatically generate inputs to test programs, and can generate inputs that cause two implementations of the same protocol to behave differently;
	\item \textbf{Automatic Network Protocol Understanding}: Dynamic taint analysis has been used to automatically understand the behavior of network protocols.
\end{itemize}

However, recalling the limitations and the problems with both analysis I would say that they have to be used carefully. Their effectiveness depends on the specific application.

	
	
	%bibliography
	\newpage
	\nocite{*}
	\printbibliography % equivale a dare un \printbibliography per ogni categoria
	
\end{document}